\documentclass[sigconf,nonacm]{acmart}

% ── Packages ──────────────────────────────────────────────────────────────────
\usepackage{listings}
\usepackage{xcolor}
\usepackage{booktabs}
\usepackage{tikz}
\usepackage{tcolorbox}
\usepackage{multicol}
\usepackage{enumitem}

\tcbuselibrary{skins,breakable}

% ── TODO box ──────────────────────────────────────────────────────────────────
\newtcolorbox{todobox}[1][]{
  colback=yellow!10,
  colframe=orange!80!black,
  fonttitle=\bfseries,
  title={TODO},
  #1
}

% ── Image placeholder box ────────────────────────────────────────────────────
\newtcolorbox{imgplaceholder}[1][]{
  colback=gray!10,
  colframe=gray!60!black,
  fonttitle=\bfseries,
  title={Figure Placeholder},
  #1
}

% ── Code listing style ───────────────────────────────────────────────────────
\definecolor{codebg}{RGB}{248,248,248}
\definecolor{keyword}{RGB}{0,0,180}
\definecolor{string}{RGB}{0,128,0}
\definecolor{comment}{RGB}{128,128,128}
\definecolor{pykey}{RGB}{175,0,219}

\lstdefinelanguage{orca}{
  keywords={model, agent, tool, task, knowledge, workflow, trigger, schema, input, edge, true, false, null},
  keywordstyle=\color{keyword}\bfseries,
  commentstyle=\color{comment}\itshape,
  stringstyle=\color{string},
  morecomment=[l]{//},
  morecomment=[l]{\#},
  morestring=[b]",
  morestring=[s]{```}{```},
  sensitive=true,
}

\lstdefinestyle{orcaStyle}{
  language=orca,
  backgroundcolor=\color{codebg},
  basicstyle=\ttfamily\scriptsize,
  breaklines=true,
  frame=single,
  framerule=0.4pt,
  rulecolor=\color{gray!50},
  numbers=left,
  numberstyle=\tiny\color{gray},
  tabsize=2,
  captionpos=b,
  aboveskip=6pt,
  belowskip=4pt,
  xleftmargin=1.5em,
  framexleftmargin=1.5em,
  showspaces=false,
  showstringspaces=false,
}

\lstdefinestyle{pythonStyle}{
  language=Python,
  backgroundcolor=\color{codebg},
  basicstyle=\ttfamily\scriptsize,
  breaklines=true,
  frame=single,
  framerule=0.4pt,
  rulecolor=\color{gray!50},
  numbers=left,
  numberstyle=\tiny\color{gray},
  tabsize=2,
  captionpos=b,
  aboveskip=6pt,
  belowskip=4pt,
  xleftmargin=1.5em,
  framexleftmargin=1.5em,
  keywordstyle=\color{pykey}\bfseries,
  stringstyle=\color{string},
  commentstyle=\color{comment}\itshape,
}

\lstset{style=orcaStyle}

% ── Metadata ─────────────────────────────────────────────────────────────────
\title{Orca: A Declarative Language for AI Agent Orchestration}

\author{Thakee Nathees}
\email{t.abdulmaj@student.curtin.edu.au}
\affiliation{%
  \institution{Curtin University}
  \department{Department of Computing}
  \city{Perth}
  \country{Australia}
}

\begin{document}

\begin{abstract}
As AI agent frameworks such as LangGraph, AutoGen, and CrewAI grow in
capability, they also grow in complexity---requiring developers to write
hundreds of lines of imperative Python to wire models, tools, prompts, and
orchestration graphs. This paper presents \textbf{Orca}, a declarative
domain-specific language (DSL) for AI agent orchestration. Orca uses an
HCL-inspired block syntax to define models, agents, tools, tasks, and
workflows in a concise, readable notation that compiles to
Python/LangGraph code. The Orca compiler performs static semantic analysis
including reference resolution, type checking, and schema validation,
catching at compile time errors that would otherwise surface only at
runtime. Generated code is fully annotated with source mappings back to
the original \texttt{.oc} source, enabling debugging at the DSL level.
We describe the language design, the four-stage compiler pipeline
(lexer, parser, analyzer, code generator), and the type system that
underpins both static checking and editor integration via the Language
Server Protocol.
\end{abstract}

\maketitle

% ══════════════════════════════════════════════════════════════════════════════
\section{Introduction}
\label{sec:intro}
% ══════════════════════════════════════════════════════════════════════════════

The rise of large language models (LLMs) has given birth to a new class of
software: \emph{AI agent systems}---programs that combine one or more LLMs
with external tools, memory, and orchestration logic to accomplish complex
tasks autonomously. Frameworks such as LangGraph~\cite{langgraph},
AutoGen~\cite{autogen}, and CrewAI~\cite{crewai} provide the primitives for
building such systems, but they require developers to express agent
configurations, tool bindings, and execution graphs in general-purpose
imperative Python code.

This approach has several drawbacks. First, \textbf{boilerplate}: defining a
simple two-agent pipeline in LangGraph requires importing modules, defining
state schemas, instantiating models, writing node functions, constructing a
state graph, adding nodes and edges, and compiling the graph---often spanning
100+ lines for what is conceptually a five-line description. Second,
\textbf{late error detection}: misconfigurations such as referencing a
non-existent tool, assigning a string where a model reference is expected, or
omitting a required field are only caught at runtime, often deep inside a
framework stack trace. Third, \textbf{poor readability}: the intent of an
agent system---which agents exist, what tools they use, and how they
orchestrate---is buried within imperative wiring code.

We present \textbf{Orca}, a domain-specific language designed to address
these problems. Orca adopts an HCL-inspired~\cite{terraform} declarative
block syntax where developers describe \emph{what} an agent system consists
of, rather than \emph{how} to wire it together. The Orca compiler translates
these declarations into fully functional Python/LangGraph code, performing
static analysis along the way.

This paper makes the following contributions:
\begin{enumerate}[leftmargin=*]
  \item A \textbf{declarative language design} for AI agent orchestration,
        with a uniform block syntax covering models, agents, tools, tasks,
        workflows, and triggers.
  \item A \textbf{four-stage compiler} (lexer, Pratt parser, semantic
        analyzer, code generator) implemented in Go, featuring error-tolerant
        parsing and a unified type system.
  \item A \textbf{static analysis framework} that catches reference errors,
        type mismatches, and schema violations at compile time, with
        diagnostic codes suitable for IDE integration.
  \item \textbf{Source-mapped code generation} that annotates every line of
        generated Python with its originating \texttt{.oc} source location,
        enabling DSL-level debugging.
\end{enumerate}

% ══════════════════════════════════════════════════════════════════════════════
\section{Background}
\label{sec:background}
% ══════════════════════════════════════════════════════════════════════════════

\subsection{AI Agent Frameworks}
\label{sec:background:frameworks}

Modern AI agent frameworks provide the building blocks for constructing
systems where LLMs interact with tools, databases, and other agents.
\textbf{LangGraph}~\cite{langgraph} models agent systems as stateful graphs
where nodes represent computation steps (agent invocations, tool calls) and
edges represent control flow, including conditional branching and cycles.
\textbf{AutoGen}~\cite{autogen} focuses on multi-agent conversation, where
agents communicate via message passing. \textbf{CrewAI}~\cite{crewai}
introduces role-based agent definitions with task assignment.

All of these frameworks share a common programming model: Python classes
and function calls. A typical LangGraph program involves:
\begin{enumerate}[leftmargin=*]
  \item Defining a \texttt{TypedDict} for the graph state.
  \item Instantiating LLM clients (\texttt{ChatOpenAI}, etc.).
  \item Writing Python functions for each graph node.
  \item Constructing a \texttt{StateGraph}, adding nodes and edges.
  \item Compiling and invoking the graph.
\end{enumerate}

This imperative approach means that the \emph{structure} of an agent
system---which agents exist, what tools they use, how they
connect---is expressed procedurally rather than declaratively, making it
difficult to extract at a glance.

\subsection{Domain-Specific Languages for Configuration}
\label{sec:background:dsls}

Domain-specific languages have a long history of success in configuration-heavy
domains. HCL (HashiCorp Configuration Language)~\cite{terraform}, used by
Terraform, allows infrastructure to be declared as blocks:

\begin{lstlisting}[language={},keywordstyle=\color{keyword}\bfseries,
  morekeywords={resource,ami,instance_type}]
resource "aws_instance" "web" {
  ami           = "ami-0c55b159"
  instance_type = "t2.micro"
}
\end{lstlisting}

This block syntax---a keyword, a name, and a body of key-value
pairs---is both human-readable and machine-parseable. It separates
\emph{what} (a resource with these properties) from \emph{how} (the
provider creates it). Orca applies this same principle to AI agent
orchestration: declare the agents, tools, and workflows; let the compiler
handle the Python wiring.

Other DSL approaches in the AI space include DSPy~\cite{dspy}, which
provides a programming model for composing LLM pipelines as optimizable
modules, and LMQL~\cite{lmql}, a query language for constrained LLM
generation. Orca differs from these in targeting the \emph{orchestration}
layer---the composition of multiple agents, tools, and workflows---rather
than individual prompt or generation control.

% ══════════════════════════════════════════════════════════════════════════════
\section{Language Design}
\label{sec:language}
% ══════════════════════════════════════════════════════════════════════════════

\subsection{Design Goals}
\label{sec:language:goals}

Orca's design is guided by five principles:

\begin{enumerate}[leftmargin=*]
  \item \textbf{Declarative over imperative.} Developers describe the
        components of an agent system and their relationships; the compiler
        handles code generation.
  \item \textbf{Conciseness.} An agent system expressible in 30 lines of
        Orca should not require 300 lines of Python.
  \item \textbf{Static safety.} The compiler catches misconfigurations---undefined
        references, type mismatches, missing required fields---before any
        code runs.
  \item \textbf{Debuggability.} Generated code carries source mappings back
        to the \texttt{.oc} file, so developers debug in terms of their
        declarations, not generated Python.
  \item \textbf{Framework independence.} The language is decoupled from any
        single backend. The current compiler targets LangGraph, but the
        design permits additional backends.
\end{enumerate}

\subsection{Syntax Overview}
\label{sec:language:syntax}

The syntax of Orca is defined by the following grammar:

\begin{definition}[Orca Grammar]\label{def:grammar}
A program $P$ is defined by the grammar:
\begin{align}
  \mathit{Program}   &::= \mathit{Block}^* \\
  \mathit{Block}     &::= \mathit{BlockType} \;\; \mathit{Ident} \;\; \texttt{\{} \;\; \mathit{Field}^* \;\; \texttt{\}} \\
  \mathit{BlockType} &::= \texttt{model} \mid \texttt{agent} \mid \texttt{tool} \mid \texttt{input} \nonumber \\
                     &\quad\mid \texttt{output} \mid \texttt{let} \mid \texttt{workflow} \nonumber \\
                     &\quad\mid \texttt{cron} \mid \texttt{event} \mid \texttt{webhook} \\
  \mathit{Field}     &::= \mathit{Ident} \;\; \texttt{=} \;\; \mathit{Expr} \\
  \mathit{Expr}      &::= \mathit{Literal} \mid \mathit{Ident} \mid \mathit{List} \mid \mathit{FlowExpr} \nonumber \\
                     &\quad\mid \mathit{Expr}\texttt{.}\mathit{Ident} \nonumber \\
                     &\quad\mid \mathit{Expr}\texttt{[}\mathit{Expr}\texttt{]} \\
  \mathit{Literal}   &::= \mathit{String} \mid \mathit{Number} \mid \mathit{Bool} \\
  \mathit{String}    &::= \texttt{"} \; (\mathit{char} \mid \texttt{\$\{} \mathit{Expr} \texttt{\}})^* \; \texttt{"} \nonumber \\
                     &\quad\mid \texttt{```} \; \mathit{char}^* \; \texttt{```} \\
  \mathit{List}      &::= \texttt{[} \; (\mathit{Expr} \; (\texttt{,} \; \mathit{Expr})^*)? \; \texttt{]} \\
  \mathit{FlowExpr}  &::= \mathit{Expr} \; (\texttt{->} \; \mathit{Expr})^+
\end{align}
\end{definition}

\noindent where $\mathit{Ident}$ denotes an identifier,
$\mathit{Expr}\texttt{.}\mathit{Ident}$ is member access (e.g.\
\texttt{builtins.web\_search}),
$\mathit{Expr}\texttt{[}\mathit{Expr}\texttt{]}$ is subscript access,
$\texttt{\$\{}\mathit{Expr}\texttt{\}}$ is string interpolation within
double-quoted strings, and \texttt{```} delimiters denote raw
multi-line strings. $\mathit{FlowExpr}$ defines agent sequencing in
workflows using the arrow operator.

Every construct is a \textbf{block}: a keyword identifying the block type,
a user-chosen name, and a brace-delimited body of key-value assignments.
Orca defines ten block types, each corresponding to a component of an
agent system. Figure~\ref{fig:full-example} shows a complete Orca program.

\begin{figure}[t]
\begin{lstlisting}[caption={A complete Orca program defining a multi-agent
  research-and-writing pipeline.},label={fig:full-example}]
model gpt4 {
  provider    = "openai"
  model_name  = "gpt-4o"
  temperature = 0.2
}

tool web_search {
  name = "web_search"
}

agent researcher {
  model   = gpt4
  tools   = [web_search]
  persona = ```
    You research topics thoroughly
    and return structured findings.
    ```
}

agent writer {
  model   = gpt4
  persona = "You are a technical writer."
}

task research {
  agent  = researcher
  prompt = "Research the given topic."
}

workflow report_pipeline {
  name = "report"
}

trigger daily {
  name     = "daily_report"
}
\end{lstlisting}
\end{figure}

\noindent\textbf{Block types.} Table~\ref{tab:blocks} summarizes each
block type and its purpose.

\begin{table}[t]
\caption{Orca block types and their roles in an agent system.}
\label{tab:blocks}
\centering
\small
\begin{tabular}{@{}ll@{}}
\toprule
\textbf{Block} & \textbf{Purpose} \\
\midrule
\texttt{model}     & LLM provider configuration (provider, model, temperature) \\
\texttt{agent}     & Agent definition (model, persona, tools) \\
\texttt{tool}      & External integration (API, builtin, function) \\
\texttt{task}      & Work unit assigned to an agent \\
\texttt{knowledge} & RAG / data source attachment \\
\texttt{workflow}  & Agent orchestration graph \\
\texttt{trigger}   & Cron, webhook, or event-based execution \\
\texttt{schema}    & User-defined type for structured data \\
\texttt{input}     & Runtime input parameter with type and default \\
\bottomrule
\end{tabular}
\end{table}

\noindent\textbf{Value types.} Assignments support string, integer, float,
boolean, and null literals; lists (\texttt{[a, b, c]}); maps
(\texttt{\{key: value\}}); and \emph{block references}---bare identifiers
that refer to other blocks by name (e.g., \texttt{model = gpt4}).

\noindent\textbf{Annotations.} Orca supports annotations prefixed with
\texttt{@}, used for metadata and compiler directives:

\begin{lstlisting}
@desc("Sampling temperature between 0 and 1.")
temperature = float | null

@suppress("unknown-field")
custom_param = "value"
\end{lstlisting}

The \texttt{@desc} annotation documents fields; \texttt{@suppress}
silences specific diagnostics when the developer intentionally deviates
from the schema.

\noindent\textbf{Workflow arrow syntax.} Workflows use an arrow operator
(\texttt{->}) to express agent execution order:

\begin{lstlisting}
workflow pipeline {
  flow = researcher -> writer -> reviewer
}
\end{lstlisting}

This compiles to a linear graph:
\texttt{START} $\rightarrow$ \texttt{researcher} $\rightarrow$
\texttt{writer} $\rightarrow$ \texttt{reviewer} $\rightarrow$ \texttt{END}.

\begin{todobox}
Add branching and parallel workflow syntax examples once implemented
in the compiler (conditional edges, parallel fan-out with
\texttt{[a, b, c] -> aggregator}).
\end{todobox}

\subsection{Type System}
\label{sec:language:types}

Orca employs a structural type system designed to validate block
configurations at compile time. The type system is built around four
type kinds:

\begin{enumerate}[leftmargin=*]
  \item \textbf{BlockRef}: Named types encompassing both primitives
        (\texttt{string}, \texttt{int}, \texttt{number}, \texttt{bool},
        \texttt{null}) and block references (\texttt{model}, \texttt{agent},
        \texttt{tool}, etc.). Every user-defined \texttt{schema} block
        also becomes a \texttt{BlockRef} type.
  \item \textbf{List}: Parameterized ordered collections, e.g.,
        \texttt{list[tool]}.
  \item \textbf{Map}: Key-value collections with string keys, e.g.,
        \texttt{map[str]}.
  \item \textbf{Union}: Disjunctive types using the pipe operator, e.g.,
        \texttt{string | model} or \texttt{number | null}.
\end{enumerate}

A key design decision is the \emph{unification of primitives and block
references}: both are represented as \texttt{BlockRef} with a
\texttt{BlockKind} string. This simplifies the type checker---there is
no special case for primitives versus user-defined types---and allows
the schema system to be extended naturally by user-defined \texttt{schema}
blocks.

\noindent\textbf{Type compatibility.} The function
\texttt{IsCompatible(got, expected)} determines whether a value of type
\texttt{got} can be assigned to a field expecting \texttt{expected}. The
rules are:
\begin{itemize}[leftmargin=*]
  \item The \texttt{any} type is compatible with every type.
  \item Unions are compatible if the \texttt{got} type matches any member.
  \item Numeric widening: \texttt{int} is compatible with \texttt{float}.
  \item Block references match by name equality.
  \item List and map types check element/value type compatibility recursively.
\end{itemize}

\noindent\textbf{Type inference.} Expression types are inferred
bottom-up from literals and symbol table lookups. For example, the
expression \texttt{[web\_search, gmail]} is inferred as
\texttt{list[tool]} if both identifiers resolve to \texttt{tool} blocks
in the symbol table.

\subsection{Block Schemas}
\label{sec:language:schemas}

Each block type has a \textbf{schema} defining its valid fields, their
types, optionality, and descriptions. Schemas serve as the contract between
the developer and the compiler: any field not in the schema is flagged,
any required field that is missing is reported, and any value whose type
does not match the schema triggers a type mismatch diagnostic.

Orca's built-in schemas are themselves defined in Orca syntax in an
embedded \texttt{builtins.oc} file, parsed at compiler initialization:

\begin{lstlisting}
schema agent {
  @desc("The model to use.")
  model = string | model

  @desc("The persona of the agent.")
  persona = string

  @desc("The tools available to this agent.")
  tools = list[tool] | null
}
\end{lstlisting}

This self-hosting approach---defining the compiler's type information in
the language it compiles---ensures that the schema notation is
battle-tested and that adding a new block type requires only adding a
schema definition, not modifying the compiler.

Users can define custom schemas using the \texttt{schema} block:

\begin{lstlisting}
schema vpc_data_t {
  region      = string
  cidr_blocks = list[string]
}

input vpc {
  type = vpc_data_t
  desc = "VPC configuration"
}
\end{lstlisting}

The compiler registers user schemas alongside built-in schemas, enabling
full type checking on \texttt{input} blocks and member access expressions
like \texttt{vpc.region}.

% ══════════════════════════════════════════════════════════════════════════════
\section{Compiler Architecture}
\label{sec:compiler}
% ══════════════════════════════════════════════════════════════════════════════

The Orca compiler is implemented in Go and follows a classical four-stage
pipeline:

\begin{center}
\texttt{.oc} $\xrightarrow{\text{Lexer}}$ Tokens
$\xrightarrow{\text{Parser}}$ AST
$\xrightarrow{\text{Analyzer}}$ Diagnostics
$\xrightarrow{\text{Codegen}}$ Python
\end{center}

\begin{imgplaceholder}
Detailed pipeline diagram showing all four stages with intermediate
representations. Left to right: source text, token stream (with example
tokens), AST tree (with example nodes), annotated AST with diagnostics,
and generated Python output. Show a small concrete example flowing through
each stage.
\end{imgplaceholder}

Each stage is implemented as an independent Go package with a clean
interface, enabling isolated testing and future extension.

\subsection{Lexer}
\label{sec:compiler:lexer}

The lexer is a hand-written, character-at-a-time scanner that produces a
stream of typed tokens. It maintains a two-character lookahead window
(\texttt{ch} and \texttt{peekChar}) and tracks line and column positions
incrementally.

Each token carries four source coordinates: \texttt{Line}, \texttt{Column},
\texttt{EndLine}, and \texttt{EndCol}. This range information is essential
for two purposes: precise diagnostic reporting (``error at line 7, columns
3--15'') and source mapping in generated code.

The lexer recognizes 39 token types organized into categories:

\begin{itemize}[leftmargin=*]
  \item \textbf{Literals}: identifiers, integers, floats, strings.
  \item \textbf{Keywords}: the nine block types plus \texttt{true},
        \texttt{false}, \texttt{null}.
  \item \textbf{Operators}: \texttt{=}, \texttt{.}, \texttt{->},
        \texttt{|}, \texttt{+}, \texttt{-}, \texttt{*}, \texttt{/},
        \texttt{@}, \texttt{?}, \texttt{:}, \texttt{\textbackslash},
        \texttt{==}, \texttt{!=}, \texttt{<}, \texttt{>},
        \texttt{<=}, \texttt{>=}.
  \item \textbf{Delimiters}: braces, brackets, parentheses, comma, colon.
\end{itemize}

\noindent\textbf{String handling.} Orca has two string types.
\emph{Double-quoted strings} (\texttt{"..."}) are single-line and support
escape sequences (\texttt{\textbackslash n}, \texttt{\textbackslash t},
\texttt{\textbackslash\textbackslash}, \texttt{\textbackslash"}).
\emph{Raw strings} use triple backticks (\texttt{```...```}) for
multi-line content with an optional language tag (e.g.
\texttt{```md}, \texttt{```py}). The closing triple backtick's column
position defines the baseline indentation, and leading whitespace up to
that column is stripped from each line. This allows developers to write:

\begin{lstlisting}
agent researcher {
  persona = ```md
    You research topics thoroughly.
    Always cite your sources.
    ```
}
\end{lstlisting}

\noindent without the indentation leaking into the string value. The
language tag enables embedded syntax highlighting in editors and is
stored on the AST node for potential use by code generation backends.

\subsection{Parser}
\label{sec:compiler:parser}

The parser combines recursive descent for top-level structure with a
\textbf{Pratt parser}~\cite{pratt1973} for expressions. It maintains a
three-token window (\texttt{prevToken}, \texttt{curToken},
\texttt{peekToken}) for one-token lookahead decisions.

\noindent\textbf{Top-level parsing.} The parser loops over the token
stream, collecting annotations and dispatching to a unified block parser
for any block keyword. All nine block types share the same syntactic
structure (\texttt{keyword name \{ assignments \}}), so a single parsing
function handles them all---the block's semantic identity is determined
by its keyword token, not by specialized grammar rules.

\noindent\textbf{Expression parsing.} The Pratt parser handles
precedence levels from lowest (\texttt{?:}) to highest (member access,
subscript, and call):

\begin{table}[t]
\caption{Operator precedence levels (lowest to highest).}
\label{tab:precedence}
\centering
\small
\begin{tabular}{@{}cll@{}}
\toprule
\textbf{Level} & \textbf{Operators} & \textbf{Associativity} \\
\midrule
1 & \texttt{?:}             & Right \\
2 & \texttt{->}             & Left \\
3 & \texttt{|}              & Left \\
4 & \texttt{==}, \texttt{!=}, \texttt{<}, \texttt{>}, \texttt{<=}, \texttt{>=} & Left \\
5 & \texttt{+}, \texttt{-}  & Left \\
6 & \texttt{*}, \texttt{/}  & Left \\
7 & \texttt{.}, \texttt{[]}, \texttt{()} & Left \\
\bottomrule
\end{tabular}
\end{table}

\noindent\textbf{Error recovery.} A critical design decision is that the
parser is \emph{error-tolerant}: it produces a (possibly partial) AST even
when the input contains syntax errors. Two recovery strategies are used:

\begin{itemize}[leftmargin=*]
  \item \texttt{syncToBlockEnd}: after a block header error, advance to the
        closing \texttt{\}} or the next block keyword.
  \item \texttt{syncToNextAssignment}: after an assignment error, advance to
        the next \texttt{identifier = ...} pattern or \texttt{\}}.
\end{itemize}

Both strategies guarantee forward progress (at least one token consumed)
to prevent infinite loops. Error-tolerant parsing is essential for
editor integration: the Language Server Protocol (LSP) server
re-parses on every keystroke, and incomplete input must not crash the
pipeline.

\noindent\textbf{Lambda expressions.} Orca supports first-class anonymous
functions via lambda expressions. The backslash token (\texttt{\textbackslash})
visually resembles the Greek letter $\lambda$ and introduces the parameter list:

\begin{lstlisting}
let funcs {
  add = \(a number, b number) number -> a + b
  double = \(x number) -> x * 2
  greet = \() -> "hello"
  add_k = \(k number) -> \(n number) -> k + n
}
\end{lstlisting}

\noindent The return type annotation after the closing parenthesis is optional;
when omitted, it is inferred from the body expression. The body is always a
single expression following the \texttt{->} arrow. Lambdas are closures---they
capture variables from enclosing scopes. The type of a lambda is
\texttt{callable[P1, \ldots, Pn, R]} where $P_i$ are parameter types and $R$
is the return type. Lambda expressions compile directly to Python
\texttt{lambda} expressions.

\noindent\textbf{Scoped symbol table.} Lambda parameters introduce a new
lexical scope. The symbol table uses a linked chain of scopes: \texttt{PushScope}
creates a child scope for lambda parameters, and \texttt{PopScope} restores the
parent. Identifier lookup walks up the chain, so inner lambdas can reference
outer parameters (closures) and outer parameters can be shadowed by inner ones.

\subsection{Semantic Analysis}
\label{sec:compiler:analyzer}

The analyzer performs two passes over the AST:

\noindent\textbf{Pass 1: Symbol table construction.} The analyzer walks
all block statements and registers each block's name in a symbol table,
mapping names to their \texttt{BlockRefType}. Built-in type names
(\texttt{string}, \texttt{number}, etc.) are pre-seeded. This
pass also detects duplicate block names.

\noindent\textbf{Pass 2: Block validation.} For each block, the analyzer
retrieves the corresponding schema and performs three categories of checks:

\begin{enumerate}[leftmargin=*]
  \item \textbf{Field coverage}: every required (non-optional) field in the
        schema must be present. Missing fields produce a
        \texttt{missing-field} diagnostic spanning the block's braces.
  \item \textbf{Field validation}: each assignment is checked against the
        schema. Unknown fields produce \texttt{unknown-field}; type
        mismatches (via \texttt{IsCompatible}) produce \texttt{type-mismatch}.
  \item \textbf{Reference resolution}: all identifier expressions are
        looked up in the symbol table. Undefined identifiers produce
        \texttt{undefined-ref}. Member access expressions (e.g.,
        \texttt{vpc.region}) are validated by resolving the object's type
        and checking the member against the type's schema.
\end{enumerate}

Table~\ref{tab:diagnostics} lists all diagnostic codes.

\begin{table}[t]
\caption{Diagnostic codes emitted by the Orca compiler.}
\label{tab:diagnostics}
\centering
\small
\begin{tabular}{@{}lll@{}}
\toprule
\textbf{Code} & \textbf{Severity} & \textbf{Description} \\
\midrule
\texttt{syntax}            & Error   & Parse error \\
\texttt{duplicate-block}   & Error   & Block name already defined \\
\texttt{duplicate-field}   & Error   & Field assigned twice in a block \\
\texttt{missing-field}     & Error   & Required schema field absent \\
\texttt{unknown-field}     & Warning & Field not defined in schema \\
\texttt{type-mismatch}     & Error   & Value incompatible with field type \\
\texttt{undefined-ref}     & Error   & Identifier not in symbol table \\
\texttt{unknown-member}    & Error   & Member not on resolved type \\
\texttt{invalid-subscript} & Error   & Non-integer index on list \\
\bottomrule
\end{tabular}
\end{table}

\noindent\textbf{Diagnostic suppression.} The \texttt{@suppress}
annotation allows developers to silence specific diagnostics:

\begin{lstlisting}
@suppress("unknown-field")
custom_param = "value"
\end{lstlisting}

Suppression can be applied at the block level (suppressing all diagnostics
within the block) or at the field level. This mechanism is essential during
language evolution: when a developer uses a field that is planned but not
yet in the schema, suppression provides an escape hatch without disabling
all checking.

\subsection{Code Generation}
\label{sec:compiler:codegen}

\begin{todobox}
The code generation backend is under active development. This section will
describe:
\begin{itemize}
  \item How each block type maps to Python/LangGraph constructs:
        \texttt{model} $\rightarrow$ \texttt{ChatOpenAI(...)},
        \texttt{agent} $\rightarrow$ node function,
        \texttt{workflow} $\rightarrow$ \texttt{StateGraph} edges.
  \item Source mapping annotations: every generated line carries a comment
        \texttt{\# source.oc:L:C} tracing back to the originating
        declaration.
  \item The generated project structure in the \texttt{build/} directory.
\end{itemize}
Show concrete input/output example once codegen is complete.
\end{todobox}

The intended generated output for the program in Figure~\ref{fig:full-example}
follows this structure:

\begin{lstlisting}[style=pythonStyle,caption={Expected generated Python
  for a two-agent workflow (illustrative).},label={fig:generated}]
# Auto-generated by Orca compiler
# Source: agents.oc

from typing import TypedDict
from langgraph.graph import StateGraph, START, END
from langchain_openai import ChatOpenAI

# --- Models ---          # agents.oc:1:1
gpt4 = ChatOpenAI(        # agents.oc:1:1
    model="gpt-4o",       # agents.oc:3:3
    temperature=0.2,       # agents.oc:4:3
)

# --- Nodes ---
def researcher_node(state):  # agents.oc:11:1
    """You research topics thoroughly..."""
    pass  # tool-bound agent invocation

def writer_node(state):      # agents.oc:20:1
    """You are a technical writer."""
    pass  # agent invocation

# --- Graph ---              # agents.oc:28:1
graph = StateGraph(dict)
graph.add_node("researcher", researcher_node)
graph.add_node("writer", writer_node)
graph.add_edge(START, "researcher")
graph.add_edge("researcher", "writer")
graph.add_edge("writer", END)

app = graph.compile()
\end{lstlisting}

% ══════════════════════════════════════════════════════════════════════════════
\section{Editor Support}
\label{sec:editor}
% ══════════════════════════════════════════════════════════════════════════════

The compiler's architecture is designed to support rich editor integration
through the Language Server Protocol (LSP). Because the parser is
error-tolerant, the LSP server can re-parse and re-analyze on every
keystroke, providing real-time feedback.

The LSP server maintains a document state for each open file, comprising
the current AST, symbol table, and diagnostics. Three features are
currently supported:

\begin{itemize}[leftmargin=*]
  \item \textbf{Diagnostics}: parse errors and semantic diagnostics are
        published as editor squiggly lines with precise ranges derived from
        token coordinates.
  \item \textbf{Hover}: hovering over a block name or field shows type
        information and \texttt{@desc} documentation from the schema.
  \item \textbf{Go-to-definition}: jumping from a block reference
        (e.g., \texttt{model = gpt4}) to the referenced block's definition.
\end{itemize}

A \textbf{cursor context resolver} determines the semantic context at any
cursor position---whether the cursor is at the top level, inside a block
body, or within a field value---and which block and schema apply. This
context drives context-aware completions: inside an \texttt{agent} block,
the completion list shows the agent schema's fields; for a \texttt{tools}
field, it offers defined \texttt{tool} blocks.

\begin{todobox}
Add screenshots of VS Code extension showing diagnostics, hover, and
completion once the extension is polished. Include the cursor context
resolution algorithm.
\end{todobox}

% ══════════════════════════════════════════════════════════════════════════════
\section{Evaluation}
\label{sec:evaluation}
% ══════════════════════════════════════════════════════════════════════════════

\subsection{Conciseness}
\label{sec:eval:conciseness}

\begin{todobox}
Quantitative comparison once codegen is complete. Plan:
\begin{itemize}
  \item Select 3--5 agent systems of increasing complexity.
  \item Implement each in both Orca and raw LangGraph Python.
  \item Compare lines of code, token count, and cyclomatic complexity.
  \item Present results in a table and bar chart.
\end{itemize}
Preliminary estimate from the illustrative example: the Orca version is
$\sim$35 lines vs.\ $\sim$120 lines of Python (3.4$\times$ reduction).
\end{todobox}

\subsection{Error Detection}
\label{sec:eval:errors}

The Orca compiler currently detects nine categories of errors at compile
time (Table~\ref{tab:diagnostics}). Each corresponds to an error that
would be a runtime exception or silent misconfiguration in Python:

\begin{itemize}[leftmargin=*]
  \item \texttt{undefined-ref}: in Python, referencing a variable before
        assignment raises \texttt{NameError} at runtime.
  \item \texttt{type-mismatch}: passing a string where a model object is
        expected may cause a cryptic \texttt{AttributeError} deep in
        LangGraph.
  \item \texttt{missing-field}: omitting a required parameter often results
        in a \texttt{TypeError} during framework initialization.
  \item \texttt{unknown-field}: misspelling a field name is silently ignored
        in dictionary-based configuration; Orca flags it immediately.
\end{itemize}

\begin{todobox}
Construct a systematic catalog: for each diagnostic code, show the Orca
error message alongside the Python exception that would occur without
static checking. Include timing comparison (compile-time instant vs.\
runtime after potentially long LLM calls).
\end{todobox}

\subsection{Source Mapping}
\label{sec:eval:sourcemap}

\begin{todobox}
Demonstrate the debugging workflow once codegen produces source-mapped
output:
\begin{enumerate}
  \item A runtime error in the generated Python stack trace.
  \item The source map annotation on the failing line.
  \item Tracing back to the originating \texttt{.oc} declaration.
\end{enumerate}
Show that developers can reason about errors in terms of their 35-line
Orca program rather than 120 lines of generated Python.
\end{todobox}

% ══════════════════════════════════════════════════════════════════════════════
\section{Related Work}
\label{sec:related}
% ══════════════════════════════════════════════════════════════════════════════

\noindent\textbf{Agent frameworks.}
LangGraph~\cite{langgraph} provides the most flexible graph-based
orchestration model but requires imperative Python for all configuration.
AutoGen~\cite{autogen} simplifies multi-agent conversation but lacks
static analysis. CrewAI~\cite{crewai} introduces role-based agents with
YAML configuration, moving toward declarative definitions but without a
type system or compiler. Orca occupies a different point in this design
space: a full language with static checking that compiles to framework
code, rather than a framework itself.

\noindent\textbf{Infrastructure DSLs.}
HCL/Terraform~\cite{terraform} is the closest analogue to Orca's approach:
a declarative block language that compiles to imperative API calls. Orca
adapts HCL's syntax philosophy---keyword-name-body blocks, references by
name, schema-driven validation---to the agent orchestration domain. Pulumi
takes the opposite approach, embedding infrastructure in general-purpose
languages; Orca argues that a dedicated syntax yields better readability
and static analysis for the agent domain.

\noindent\textbf{AI programming DSLs.}
DSPy~\cite{dspy} provides a programming model for composing and optimizing
LLM pipelines as declarative modules. LMQL~\cite{lmql} extends Python with
a query syntax for constrained LLM generation. Both target individual
model interactions rather than multi-agent orchestration. Orca is
complementary: it could potentially use DSPy modules as tool
implementations within an Orca-orchestrated system.

% ══════════════════════════════════════════════════════════════════════════════
\section{Future Work}
\label{sec:future}
% ══════════════════════════════════════════════════════════════════════════════

Several extensions are planned:

\begin{itemize}[leftmargin=*]
  \item \textbf{Additional codegen backends.} The language-backend separation
        allows targeting frameworks beyond LangGraph, such as AutoGen or
        direct API calls.
  \item \textbf{Visual editor.} A node-and-edge graph editor that generates
        \texttt{.oc} files, providing a visual frontend for users who prefer
        graphical workflow design.
  \item \textbf{Full LSP server.} Expanding editor support with rename
        refactoring, find-all-references, code actions, and workspace-wide
        diagnostics.
  \item \textbf{Runtime debugging.} Integrating source maps with a debugger
        that allows breakpoints and stepping in terms of \texttt{.oc}
        source lines.
  \item \textbf{Package system.} A module system for reusable agent
        components (e.g., a pre-built \texttt{web\_researcher} agent
        importable across projects).
\end{itemize}

% ══════════════════════════════════════════════════════════════════════════════
\section{Conclusion}
\label{sec:conclusion}
% ══════════════════════════════════════════════════════════════════════════════

We have presented Orca, a declarative domain-specific language for AI
agent orchestration. Orca's HCL-inspired block syntax allows developers
to define models, agents, tools, tasks, and workflows concisely, while
the compiler's static analysis catches misconfigurations at compile time
that would otherwise surface as runtime errors. The four-stage compiler
pipeline---lexer, Pratt parser with error recovery, semantic analyzer with
a unified type system, and source-mapped code generator---demonstrates that
the agent orchestration domain is amenable to the same
declare-then-compile approach that has proven successful for infrastructure
(Terraform/HCL) and data pipelines.

As AI agent systems grow in complexity, the gap between declarative intent
and imperative implementation widens. Orca aims to bridge this gap:
developers describe \emph{what} their agent system does, and the compiler
handles \emph{how}. We believe this \emph{agents-as-code} paradigm,
analogous to infrastructure-as-code, will become increasingly important as
agent systems move from prototypes to production.

% ══════════════════════════════════════════════════════════════════════════════
% References
% ══════════════════════════════════════════════════════════════════════════════
\bibliographystyle{ACM-Reference-Format}
\bibliography{references}

\end{document}
